% Options for packages loaded elsewhere
\PassOptionsToPackage{unicode}{hyperref}
\PassOptionsToPackage{hyphens}{url}
\documentclass[
  12pt,
]{ctexart}
\usepackage{xcolor}
\usepackage{amsmath,amssymb}
\setcounter{secnumdepth}{-\maxdimen} % remove section numbering
\usepackage{iftex}
\ifPDFTeX
  \usepackage[T1]{fontenc}
  \usepackage[utf8]{inputenc}
  \usepackage{textcomp} % provide euro and other symbols
\else % if luatex or xetex
  \usepackage{unicode-math} % this also loads fontspec
  \defaultfontfeatures{Scale=MatchLowercase}
  \defaultfontfeatures[\rmfamily]{Ligatures=TeX,Scale=1}
\fi
\usepackage{lmodern}
\ifPDFTeX\else
  % xetex/luatex font selection
\fi
% Use upquote if available, for straight quotes in verbatim environments
\IfFileExists{upquote.sty}{\usepackage{upquote}}{}
\IfFileExists{microtype.sty}{% use microtype if available
  \usepackage[]{microtype}
  \UseMicrotypeSet[protrusion]{basicmath} % disable protrusion for tt fonts
}{}
\makeatletter
\@ifundefined{KOMAClassName}{% if non-KOMA class
  \IfFileExists{parskip.sty}{%
    \usepackage{parskip}
  }{% else
    \setlength{\parindent}{0pt}
    \setlength{\parskip}{6pt plus 2pt minus 1pt}}
}{% if KOMA class
  \KOMAoptions{parskip=half}}
\makeatother
\usepackage{longtable,booktabs,array}
\usepackage{calc} % for calculating minipage widths
% Correct order of tables after \paragraph or \subparagraph
\usepackage{etoolbox}
\makeatletter
\patchcmd\longtable{\par}{\if@noskipsec\mbox{}\fi\par}{}{}
\makeatother
% Allow footnotes in longtable head/foot
\IfFileExists{footnotehyper.sty}{\usepackage{footnotehyper}}{\usepackage{footnote}}
\makesavenoteenv{longtable}
\setlength{\emergencystretch}{3em} % prevent overfull lines
\providecommand{\tightlist}{%
  \setlength{\itemsep}{0pt}\setlength{\parskip}{0pt}}
\usepackage{bookmark}
\IfFileExists{xurl.sty}{\usepackage{xurl}}{} % add URL line breaks if available
\urlstyle{same}
\hypersetup{
  hidelinks,
  pdfcreator={LaTeX via pandoc}}

\author{SCX}
\date{2026}

\begin{document}

\section{SCX 终极 TODO 列表}\label{scx-ux7ec8ux6781-todo-ux5217ux8868}

\begin{quote}
生成时间：2026-06-26 来源：综合 GPT 讨论
(\texttt{与GPT的讨论2026-06-26-07-49.md}，9011行) 目标：一份可执行的终极
TODO 列表，标注优先级、依赖、预估时间、维度和可立即执行项
\end{quote}

\begin{center}\rule{0.5\linewidth}{0.5pt}\end{center}

\subsection{2026-06-26 执行状态}\label{ux6267ux884cux72b6ux6001}

\begin{quote}
本文档的 P0/P1 项已有大量完成。完成项已标记 ✅。详细的开发待办已迁移到
\texttt{SCX\_TODO\_2026-06-26.md}。
\end{quote}

\subsubsection{Phase 0 完成项}\label{phase-0-ux5b8cux6210ux9879}

\begin{longtable}[]{@{}
  >{\raggedright\arraybackslash}p{(\linewidth - 6\tabcolsep) * \real{0.1429}}
  >{\raggedright\arraybackslash}p{(\linewidth - 6\tabcolsep) * \real{0.2857}}
  >{\raggedright\arraybackslash}p{(\linewidth - 6\tabcolsep) * \real{0.2857}}
  >{\raggedright\arraybackslash}p{(\linewidth - 6\tabcolsep) * \real{0.2857}}@{}}
\toprule\noalign{}
\begin{minipage}[b]{\linewidth}\raggedright
\#
\end{minipage} & \begin{minipage}[b]{\linewidth}\raggedright
TODO
\end{minipage} & \begin{minipage}[b]{\linewidth}\raggedright
状态
\end{minipage} & \begin{minipage}[b]{\linewidth}\raggedright
备注
\end{minipage} \\
\midrule\noalign{}
\endhead
\bottomrule\noalign{}
\endlastfoot
0.1 & 私人 SCX 仓库 & ✅ & 已建立于
\texttt{G:/Xiaogan\_Supercomputing\_data/SCX/} \\
0.2 & \texttt{DEVELOPMENT\_LOG.md} & ✅ & + \texttt{IP\_NOTE.md} +
\texttt{CLEAN\_ROOM\_CHECK.md} 三文件完成 \\
0.3 & 代码来源整理 & ✅ & Clean-room 检查：30 源文件零学校引用 \\
0.4 & 协议检查 & ✅ & 确认无签署任何 IP 相关协议 \\
0.5 & Clean-room 重写 & ✅ & PASSED --- SCX v0.3.0 全 clean-room \\
0.6 & 成果归属说明 & ✅ & \texttt{IP\_NOTE.md} 完成 \\
0.7 & 不上传学校服务器 & ✅ & SCX 核心代码仅存个人电脑 \\
0.8 & 不用学校邮箱 & ✅ & 已处理 \\
0.9 & EGP Paper 1 草稿 & ◐ 部分完成 & Introduction + Methods
框架已写，等 DFT 结果补实验 \\
0.10 & SCX-Theory 草稿 & ◐ 部分完成 & 定义 + 命题完成，等合成实验 Figure
后发 arXiv \\
0.11 & 整理 v0.3.0 & ✅ & 35 Python 文件, 336 tests, 全部通过 \\
0.12 & Implementation note & ✅ & SCX 骨架文档化 \\
0.13 & ACE pipeline 整理 & ✅ & 代码 clean-room 重写完成 \\
0.14 & VASP 输入检查 & ✅ & GaN\_v1 (540 jobs) + AlN\_v4 (164 jobs) \\
0.15 & VASP 任务提交 & ✅ & 已提交超算运行中 \\
0.16 & 监控脚本 & ✅ & 已配置 \\
\end{longtable}

\subsubsection{Phase 1 完成项}\label{phase-1-ux5b8cux6210ux9879}

\begin{longtable}[]{@{}llll@{}}
\toprule\noalign{}
\# & TODO & 状态 & 备注 \\
\midrule\noalign{}
\endhead
\bottomrule\noalign{}
\endlastfoot
1.9 & SCX-Health 仓库 & ✅ & \texttt{scx-health/} Apache 2.0 开源准备 \\
1.10 & MedMNIST Compress 实验 & ✅ & 30\% 压缩 +6.00\% 精度 \\
1.11 & MedMNIST Noise 实验 & ✅ & SCX-Noise F1=0.617 \\
1.12 & MedMNIST Routing 实验 & ✅ & 多弱专家路由验证 \\
1.15 & SCX one-pager & ✅ & 策略文档已包含 \\
1.16 & 开放策略文档 & ✅ & 三层开放策略已定义 \\
1.4 & 时间戳存证 & ✅ & Sha-256 + 开发日志持续记录 \\
\end{longtable}

\subsubsection{等待中的项}\label{ux7b49ux5f85ux4e2dux7684ux9879}

\begin{longtable}[]{@{}llll@{}}
\toprule\noalign{}
\# & TODO & 阻塞因素 & 预计解锁 \\
\midrule\noalign{}
\endhead
\bottomrule\noalign{}
\endlastfoot
1.5 & SCX-MLIP 论文 & DFT 结果未到 & 2-4 周 \\
1.6 & 投稿 & 1.5 未完成 & 2-4 周 \\
1.7 & arXiv 占坑 & 等待合成实验 Figure & 立即可做 \\
1.13 & HAM10000/ISIC 实验 & 等待 GPU & 8 月中旬 \\
1.14 & SCX-Health 论文 & 等待 GPU 重跑 & 8 月中旬 \\
1.17 & 华为合作调研 & 待启动 & 立即可做 \\
1.19-1.21 & LLM 方向 & 非主战场，低优先级 & --- \\
\end{longtable}

\begin{center}\rule{0.5\linewidth}{0.5pt}\end{center}

\subsection{维度和标记说明}\label{ux7ef4ux5ea6ux548cux6807ux8bb0ux8bf4ux660e}

\begin{longtable}[]{@{}lll@{}}
\toprule\noalign{}
标记 & 维度 & 说明 \\
\midrule\noalign{}
\endhead
\bottomrule\noalign{}
\endlastfoot
{[}A{]} & Academic / 学术 & 论文、投稿、学术影响力 \\
{[}C{]} & Code / 代码 & 软件工程、实现、开源 \\
{[}B{]} & Business / 商业 & 商业模式、合作、收入 \\
{[}IP{]} & Intellectual Property / 知识产权 & 专利、权属、保护 \\
{[}LLM{]} & LLM / 大模型 & 大模型数据评价方向 \\
\end{longtable}

\textbf{优先级}: P0=立即做, P1=1-2月内, P2=3-6月内, P3=6-12月

\textbf{依赖标记}: \texttt{A\ →\ B} 表示 A 阻塞 B

\textbf{时间标记}: ⚡=现在就能做, ⏳=需要等待前置条件

\begin{center}\rule{0.5\linewidth}{0.5pt}\end{center}

\subsection{Phase 0:
立即（本周，零成本）}\label{phase-0-ux7acbux5373ux672cux5468ux96f6ux6210ux672c}

\subsubsection{IP 权属与证据链 ------
第一优先级}\label{ip-ux6743ux5c5eux4e0eux8bc1ux636eux94fe-ux7b2cux4e00ux4f18ux5148ux7ea7}

\begin{longtable}[]{@{}
  >{\raggedright\arraybackslash}p{(\linewidth - 12\tabcolsep) * \real{0.0208}}
  >{\raggedright\arraybackslash}p{(\linewidth - 12\tabcolsep) * \real{0.5903}}
  >{\raggedright\arraybackslash}p{(\linewidth - 12\tabcolsep) * \real{0.2292}}
  >{\raggedright\arraybackslash}p{(\linewidth - 12\tabcolsep) * \real{0.0278}}
  >{\raggedright\arraybackslash}p{(\linewidth - 12\tabcolsep) * \real{0.0278}}
  >{\raggedright\arraybackslash}p{(\linewidth - 12\tabcolsep) * \real{0.0486}}
  >{\raggedright\arraybackslash}p{(\linewidth - 12\tabcolsep) * \real{0.0556}}@{}}
\toprule\noalign{}
\begin{minipage}[b]{\linewidth}\raggedright
\#
\end{minipage} & \begin{minipage}[b]{\linewidth}\raggedright
TODO
\end{minipage} & \begin{minipage}[b]{\linewidth}\raggedright
优先级
\end{minipage} & \begin{minipage}[b]{\linewidth}\raggedright
依赖
\end{minipage} & \begin{minipage}[b]{\linewidth}\raggedright
预估
\end{minipage} & \begin{minipage}[b]{\linewidth}\raggedright
维度
\end{minipage} & \begin{minipage}[b]{\linewidth}\raggedright
可执行性
\end{minipage} \\
\midrule\noalign{}
\endhead
\bottomrule\noalign{}
\endlastfoot
0.1 & 新建私人 SCX 仓库（个人邮箱、个人
GitHub/GitLab，个人电脑），不要再放学校资源 & P0已完成 & 无 & 1h &
{[}IP{]}{[}C{]} & ⚡ \\
0.2 & 写 \texttt{DEVELOPMENT\_LOG.md}，记录 SCX
核心思想的产生时间、开发地点、设备、数据来源 & P0 ✅ & 无 & 2h &
{[}IP{]} & ⚡ \\
0.3 &
把所有旧代码/旧数据的来源整理成表，标注哪些用了学校资源、哪些是个人独立产生
& P0 ✅ & 无 & 3h & {[}IP{]} & ⚡ \\
0.4 & 查你签过的所有协议：入学协议、助研/奖学金协议、课题组规定、学校 IP
政策、超算使用协议 & P0 ✅ & 无 & 4h & {[}IP{]} & ⚡ \\
0.5 & SCX 代码 clean-room
重写：不要从旧项目复制粘贴，数学思想保留但代码/数据/注释全部重写 & P0 ✅
& 0.1 & 8h & {[}IP{]}{[}C{]} & ⚡ \\
0.6 & 写私人''成果归属说明''文档（SCX 项目独立于 AlGaN/DFT
的证据），存档但不对外 & P0 ✅ & 无 & 1h & {[}IP{]} & ⚡ \\
0.7 & 停止将 SCX 核心代码上传学校服务器/使用学校超算跑 SCX 实验 & P0 ✅
& 无 & 0h & {[}IP{]} & ⚡ \\
0.8 & 不再用学校邮箱注册 SCX 相关仓库，repo 权限不给课题组任何人 & P0 ✅
已处理 & 无 & 0h & {[}IP{]} & ⚡ \\
\end{longtable}

\subsubsection{论文策略 ------
快速占位}\label{ux8bbaux6587ux7b56ux7565-ux5febux901fux5360ux4f4d}

\begin{longtable}[]{@{}
  >{\raggedright\arraybackslash}p{(\linewidth - 12\tabcolsep) * \real{0.0286}}
  >{\raggedright\arraybackslash}p{(\linewidth - 12\tabcolsep) * \real{0.6643}}
  >{\raggedright\arraybackslash}p{(\linewidth - 12\tabcolsep) * \real{0.0429}}
  >{\raggedright\arraybackslash}p{(\linewidth - 12\tabcolsep) * \real{0.1357}}
  >{\raggedright\arraybackslash}p{(\linewidth - 12\tabcolsep) * \real{0.0286}}
  >{\raggedright\arraybackslash}p{(\linewidth - 12\tabcolsep) * \real{0.0429}}
  >{\raggedright\arraybackslash}p{(\linewidth - 12\tabcolsep) * \real{0.0571}}@{}}
\toprule\noalign{}
\begin{minipage}[b]{\linewidth}\raggedright
\#
\end{minipage} & \begin{minipage}[b]{\linewidth}\raggedright
TODO
\end{minipage} & \begin{minipage}[b]{\linewidth}\raggedright
优先级
\end{minipage} & \begin{minipage}[b]{\linewidth}\raggedright
依赖
\end{minipage} & \begin{minipage}[b]{\linewidth}\raggedright
预估
\end{minipage} & \begin{minipage}[b]{\linewidth}\raggedright
维度
\end{minipage} & \begin{minipage}[b]{\linewidth}\raggedright
可执行性
\end{minipage} \\
\midrule\noalign{}
\endhead
\bottomrule\noalign{}
\endlastfoot
0.9 & 写 EGP Paper 1 草稿（ACE gauge merging、energy alignment、expert
merging） & P0 & 无（已有 ACE 资产 + AlN v3 DFT 数据） & 40h &
{[}A{]}{[}C{]} & ⚡ \\
0.10 & 写 SCX-Theory 论文草稿（arXiv
先占坑）：定义、命题、可识别性边界、合成实验 & P0 & 0.5 & 30h & {[}A{]}
& ⚡ \\
\end{longtable}

\subsubsection{代码基础}\label{ux4ee3ux7801ux57faux7840}

\begin{longtable}[]{@{}
  >{\raggedright\arraybackslash}p{(\linewidth - 12\tabcolsep) * \real{0.0274}}
  >{\raggedright\arraybackslash}p{(\linewidth - 12\tabcolsep) * \real{0.6849}}
  >{\raggedright\arraybackslash}p{(\linewidth - 12\tabcolsep) * \real{0.1096}}
  >{\raggedright\arraybackslash}p{(\linewidth - 12\tabcolsep) * \real{0.0548}}
  >{\raggedright\arraybackslash}p{(\linewidth - 12\tabcolsep) * \real{0.0274}}
  >{\raggedright\arraybackslash}p{(\linewidth - 12\tabcolsep) * \real{0.0411}}
  >{\raggedright\arraybackslash}p{(\linewidth - 12\tabcolsep) * \real{0.0548}}@{}}
\toprule\noalign{}
\begin{minipage}[b]{\linewidth}\raggedright
\#
\end{minipage} & \begin{minipage}[b]{\linewidth}\raggedright
TODO
\end{minipage} & \begin{minipage}[b]{\linewidth}\raggedright
优先级
\end{minipage} & \begin{minipage}[b]{\linewidth}\raggedright
依赖
\end{minipage} & \begin{minipage}[b]{\linewidth}\raggedright
预估
\end{minipage} & \begin{minipage}[b]{\linewidth}\raggedright
维度
\end{minipage} & \begin{minipage}[b]{\linewidth}\raggedright
可执行性
\end{minipage} \\
\midrule\noalign{}
\endhead
\bottomrule\noalign{}
\endlastfoot
0.11 & 整理已有 SCX Python 包 (v0.3.0) 的 clean-room
版本，确认没有混入学校代码 & P0 ✅ & 0.1, 0.5 & 6h & {[}C{]} & ⚡ \\
0.12 & 写 SCX implementation note (method
fix)：Input→Residual→State→Statistics→Action→Validation 骨架 & P0 ✅ &
无 & 3h & {[}A{]}{[}C{]} & ⚡ \\
0.13 & 将 ACE descriptor + residual map + redundancy pipeline
整理成可复现脚本 & P0 ✅ & 0.11 & 8h & {[}C{]}{[}A{]} & ⚡ \\
\end{longtable}

\subsubsection{VASP 任务 ------
超算计算}\label{vasp-ux4efbux52a1-ux8d85ux7b97ux8ba1ux7b97}

\begin{longtable}[]{@{}
  >{\raggedright\arraybackslash}p{(\linewidth - 12\tabcolsep) * \real{0.0388}}
  >{\raggedright\arraybackslash}p{(\linewidth - 12\tabcolsep) * \real{0.6505}}
  >{\raggedright\arraybackslash}p{(\linewidth - 12\tabcolsep) * \real{0.0971}}
  >{\raggedright\arraybackslash}p{(\linewidth - 12\tabcolsep) * \real{0.0388}}
  >{\raggedright\arraybackslash}p{(\linewidth - 12\tabcolsep) * \real{0.0388}}
  >{\raggedright\arraybackslash}p{(\linewidth - 12\tabcolsep) * \real{0.0583}}
  >{\raggedright\arraybackslash}p{(\linewidth - 12\tabcolsep) * \real{0.0777}}@{}}
\toprule\noalign{}
\begin{minipage}[b]{\linewidth}\raggedright
\#
\end{minipage} & \begin{minipage}[b]{\linewidth}\raggedright
TODO
\end{minipage} & \begin{minipage}[b]{\linewidth}\raggedright
优先级
\end{minipage} & \begin{minipage}[b]{\linewidth}\raggedright
依赖
\end{minipage} & \begin{minipage}[b]{\linewidth}\raggedright
预估
\end{minipage} & \begin{minipage}[b]{\linewidth}\raggedright
维度
\end{minipage} & \begin{minipage}[b]{\linewidth}\raggedright
可执行性
\end{minipage} \\
\midrule\noalign{}
\endhead
\bottomrule\noalign{}
\endlastfoot
0.14 & 检查 GaN\_v1（540 jobs）和 AlN\_v4\_targeted（164
jobs）输入文件完整性 & P0 ✅ 已处理 & 无 & 2h & {[}A{]}{[}C{]} & ⚡ \\
0.15 & 分批提交 VASP 任务（注意 240核负载，不要一次占满） & P0 ✅ 已提交
& 0.14 & 4h & {[}A{]}{[}C{]} & ⚡ \\
0.16 & 设置任务监控脚本，自动化检查收敛状态（grep OUTCAR energy） & P0
✅ 已处理 & 0.15 & 2h & {[}C{]} & ⚡ \\
\end{longtable}

\begin{center}\rule{0.5\linewidth}{0.5pt}\end{center}

\subsection{Phase 1: 短期（1-2
月）}\label{phase-1-ux77edux671f1-2-ux6708}

\subsubsection{IP 与权属}\label{ip-ux4e0eux6743ux5c5e}

\begin{longtable}[]{@{}
  >{\raggedright\arraybackslash}p{(\linewidth - 12\tabcolsep) * \real{0.0169}}
  >{\raggedright\arraybackslash}p{(\linewidth - 12\tabcolsep) * \real{0.4576}}
  >{\raggedright\arraybackslash}p{(\linewidth - 12\tabcolsep) * \real{0.3842}}
  >{\raggedright\arraybackslash}p{(\linewidth - 12\tabcolsep) * \real{0.0226}}
  >{\raggedright\arraybackslash}p{(\linewidth - 12\tabcolsep) * \real{0.0339}}
  >{\raggedright\arraybackslash}p{(\linewidth - 12\tabcolsep) * \real{0.0395}}
  >{\raggedright\arraybackslash}p{(\linewidth - 12\tabcolsep) * \real{0.0452}}@{}}
\toprule\noalign{}
\begin{minipage}[b]{\linewidth}\raggedright
\#
\end{minipage} & \begin{minipage}[b]{\linewidth}\raggedright
TODO
\end{minipage} & \begin{minipage}[b]{\linewidth}\raggedright
优先级
\end{minipage} & \begin{minipage}[b]{\linewidth}\raggedright
依赖
\end{minipage} & \begin{minipage}[b]{\linewidth}\raggedright
预估
\end{minipage} & \begin{minipage}[b]{\linewidth}\raggedright
维度
\end{minipage} & \begin{minipage}[b]{\linewidth}\raggedright
可执行性
\end{minipage} \\
\midrule\noalign{}
\endhead
\bottomrule\noalign{}
\endlastfoot
1.1 & 找校外 IP 律师/专利代理人做一次正式权属评估 & P1 & 0.4 & 1-2 周 &
{[}IP{]} & ⚡ \\
1.2 & 判断 SCX
核心流程是否有专利价值（数据压缩、专家路由、状态价值函数、审计报告系统）
& P1 & 1.1 & 附带 & {[}IP{]} & ⏳等 1.1 \\
1.3 & 如值得，先提交发明披露/专利预审，再发 arXiv & P1 & 1.2 & 1-3 月 &
{[}IP{]}{[}A{]} & ⏳等 1.2 \\
1.4 & 关键代码做时间戳存证（Sha-256 + 可信第三方/区块链存证） & P1 ✅
已处理 & 0.1 & 2h & {[}IP{]} & ⚡ \\
\end{longtable}

\subsubsection{论文发表}\label{ux8bbaux6587ux53d1ux8868}

\begin{longtable}[]{@{}
  >{\raggedright\arraybackslash}p{(\linewidth - 12\tabcolsep) * \real{0.0196}}
  >{\raggedright\arraybackslash}p{(\linewidth - 12\tabcolsep) * \real{0.6536}}
  >{\raggedright\arraybackslash}p{(\linewidth - 12\tabcolsep) * \real{0.0392}}
  >{\raggedright\arraybackslash}p{(\linewidth - 12\tabcolsep) * \real{0.1307}}
  >{\raggedright\arraybackslash}p{(\linewidth - 12\tabcolsep) * \real{0.0392}}
  >{\raggedright\arraybackslash}p{(\linewidth - 12\tabcolsep) * \real{0.0261}}
  >{\raggedright\arraybackslash}p{(\linewidth - 12\tabcolsep) * \real{0.0915}}@{}}
\toprule\noalign{}
\begin{minipage}[b]{\linewidth}\raggedright
\#
\end{minipage} & \begin{minipage}[b]{\linewidth}\raggedright
TODO
\end{minipage} & \begin{minipage}[b]{\linewidth}\raggedright
优先级
\end{minipage} & \begin{minipage}[b]{\linewidth}\raggedright
依赖
\end{minipage} & \begin{minipage}[b]{\linewidth}\raggedright
预估
\end{minipage} & \begin{minipage}[b]{\linewidth}\raggedright
维度
\end{minipage} & \begin{minipage}[b]{\linewidth}\raggedright
可执行性
\end{minipage} \\
\midrule\noalign{}
\endhead
\bottomrule\noalign{}
\endlastfoot
1.5 & 完成 SCX-MLIP 论文（gauge-normalized expert potentials +
residual-state maps） & P1 & 0.9, 0.13, VASP 结果 & 60h & {[}A{]} & ⏳等
VASP 结果 \\
1.6 & 投稿目标：npj Computational Materials / Nature Communications /
Computer Physics Communications & P1 & 1.5 & 投稿日 & {[}A{]} & ⏳ \\
1.7 & SCX-Theory 发 arXiv 占概念（State-Conditioned eXpertise: From
Expert Reliability to Data Valuation） & P1 & 0.10 & 2 周 & {[}A{]} &
⚡ \\
1.8 & SCX-Theory 投稿 TMLR / AISTATS / SIMODS / JMLR & P2 & 1.7 & 1-2 月
& {[}A{]} & ⏳ \\
\end{longtable}

\subsubsection{医学数据实验（开源 +
论文双用）你可以现在给我建立一个专用的文件夹来存放医学数据和代码，也可以调取一个agent去下载数据}\label{ux533bux5b66ux6570ux636eux5b9eux9a8cux5f00ux6e90-ux8bbaux6587ux53ccux7528ux4f60ux53efux4ee5ux73b0ux5728ux7ed9ux6211ux5efaux7acbux4e00ux4e2aux4e13ux7528ux7684ux6587ux4ef6ux5939ux6765ux5b58ux653eux533bux5b66ux6570ux636eux548cux4ee3ux7801ux4e5fux53efux4ee5ux8c03ux53d6ux4e00ux4e2aagentux53bbux4e0bux8f7dux6570ux636e}

\begin{longtable}[]{@{}
  >{\raggedright\arraybackslash}p{(\linewidth - 12\tabcolsep) * \real{0.0328}}
  >{\raggedright\arraybackslash}p{(\linewidth - 12\tabcolsep) * \real{0.6967}}
  >{\raggedright\arraybackslash}p{(\linewidth - 12\tabcolsep) * \real{0.0492}}
  >{\raggedright\arraybackslash}p{(\linewidth - 12\tabcolsep) * \real{0.0738}}
  >{\raggedright\arraybackslash}p{(\linewidth - 12\tabcolsep) * \real{0.0328}}
  >{\raggedright\arraybackslash}p{(\linewidth - 12\tabcolsep) * \real{0.0492}}
  >{\raggedright\arraybackslash}p{(\linewidth - 12\tabcolsep) * \real{0.0656}}@{}}
\toprule\noalign{}
\begin{minipage}[b]{\linewidth}\raggedright
\#
\end{minipage} & \begin{minipage}[b]{\linewidth}\raggedright
TODO
\end{minipage} & \begin{minipage}[b]{\linewidth}\raggedright
优先级
\end{minipage} & \begin{minipage}[b]{\linewidth}\raggedright
依赖
\end{minipage} & \begin{minipage}[b]{\linewidth}\raggedright
预估
\end{minipage} & \begin{minipage}[b]{\linewidth}\raggedright
维度
\end{minipage} & \begin{minipage}[b]{\linewidth}\raggedright
可执行性
\end{minipage} \\
\midrule\noalign{}
\endhead
\bottomrule\noalign{}
\endlastfoot
1.9 & 搭建 SCX-Health 开源仓库（独立于主仓库，公开） & P1 ✅ & 0.11 & 4h
& {[}C{]}{[}A{]} & ⚡ \\
1.10 & MedMNIST v2 实验：SCX-Compress（冗余压缩 20\%-40\% 保持性能） &
P1 ✅ & 1.9 & 10h & {[}C{]}{[}A{]} & ⚡ \\
1.11 & MedMNIST v2 实验：SCX-Noise（人工标签噪声，区分 noisy state 和
learnable high-error） & P1 ✅ & 1.9 & 10h & {[}C{]}{[}A{]} & ⚡ \\
1.12 & MedMNIST v2 实验：SCX-Routing（多弱专家，状态条件路由优于平均
ensemble） & P1 ✅ & 1.9 & 10h & {[}C{]}{[}A{]} & ⚡ \\
1.13 & HAM10000/ISIC 实验：类别长尾、状态异质性、冗余压缩 & P1 & 1.9 &
15h & {[}C{]}{[}A{]} & ⚡ \\
1.14 & 写 SCX-Health 论文草稿（目标 npj Digital Medicine / Nature
Communications） & P2 & 1.10-1.13 & 40h & {[}A{]} & ⏳ \\
\end{longtable}

\subsubsection{商业基础（零成本准备）慢慢推进，等我显卡买回来再说，我8月中旬才能从合肥出差回来}\label{ux5546ux4e1aux57faux7840ux96f6ux6210ux672cux51c6ux5907ux6162ux6162ux63a8ux8fdbux7b49ux6211ux663eux5361ux4e70ux56deux6765ux518dux8bf4ux62118ux6708ux4e2dux65ecux624dux80fdux4eceux5408ux80a5ux51faux5deeux56deux6765}

\begin{longtable}[]{@{}
  >{\raggedright\arraybackslash}p{(\linewidth - 12\tabcolsep) * \real{0.0370}}
  >{\raggedright\arraybackslash}p{(\linewidth - 12\tabcolsep) * \real{0.6944}}
  >{\raggedright\arraybackslash}p{(\linewidth - 12\tabcolsep) * \real{0.0556}}
  >{\raggedright\arraybackslash}p{(\linewidth - 12\tabcolsep) * \real{0.0370}}
  >{\raggedright\arraybackslash}p{(\linewidth - 12\tabcolsep) * \real{0.0370}}
  >{\raggedright\arraybackslash}p{(\linewidth - 12\tabcolsep) * \real{0.0648}}
  >{\raggedright\arraybackslash}p{(\linewidth - 12\tabcolsep) * \real{0.0741}}@{}}
\toprule\noalign{}
\begin{minipage}[b]{\linewidth}\raggedright
\#
\end{minipage} & \begin{minipage}[b]{\linewidth}\raggedright
TODO
\end{minipage} & \begin{minipage}[b]{\linewidth}\raggedright
优先级
\end{minipage} & \begin{minipage}[b]{\linewidth}\raggedright
依赖
\end{minipage} & \begin{minipage}[b]{\linewidth}\raggedright
预估
\end{minipage} & \begin{minipage}[b]{\linewidth}\raggedright
维度
\end{minipage} & \begin{minipage}[b]{\linewidth}\raggedright
可执行性
\end{minipage} \\
\midrule\noalign{}
\endhead
\bottomrule\noalign{}
\endlastfoot
1.15 & 写 SCX one-pager
非保密版（痛点、方法、demo、收益、合作模式，不列核心实现） & P1 ✅ & 无
& 6h & {[}B{]} & ⚡ \\
1.16 & 定义三层开放策略文档：论文公开级 / 开源 reference 级 / 商业闭源级
& P1 ✅ & 无 & 3h & {[}B{]}{[}IP{]} & ⚡ \\
1.17 & 调研华为健康/穿戴团队合作入口（是否有公开合作通道/open innovation
lab） & P1 & 1.9 & 4h & {[}B{]} & ⚡ \\
1.18 & 调研国产 EDA/TCAD 公司与材料计算公司的潜在技术合作 & P2 & 1.15 &
4h & {[}B{]} & ⚡ \\
\end{longtable}

\subsubsection{大模型方向（最小验证）}\label{ux5927ux6a21ux578bux65b9ux5411ux6700ux5c0fux9a8cux8bc1}

\begin{longtable}[]{@{}
  >{\raggedright\arraybackslash}p{(\linewidth - 12\tabcolsep) * \real{0.0351}}
  >{\raggedright\arraybackslash}p{(\linewidth - 12\tabcolsep) * \real{0.7018}}
  >{\raggedright\arraybackslash}p{(\linewidth - 12\tabcolsep) * \real{0.0526}}
  >{\raggedright\arraybackslash}p{(\linewidth - 12\tabcolsep) * \real{0.0351}}
  >{\raggedright\arraybackslash}p{(\linewidth - 12\tabcolsep) * \real{0.0351}}
  >{\raggedright\arraybackslash}p{(\linewidth - 12\tabcolsep) * \real{0.0702}}
  >{\raggedright\arraybackslash}p{(\linewidth - 12\tabcolsep) * \real{0.0702}}@{}}
\toprule\noalign{}
\begin{minipage}[b]{\linewidth}\raggedright
\#
\end{minipage} & \begin{minipage}[b]{\linewidth}\raggedright
TODO
\end{minipage} & \begin{minipage}[b]{\linewidth}\raggedright
优先级
\end{minipage} & \begin{minipage}[b]{\linewidth}\raggedright
依赖
\end{minipage} & \begin{minipage}[b]{\linewidth}\raggedright
预估
\end{minipage} & \begin{minipage}[b]{\linewidth}\raggedright
维度
\end{minipage} & \begin{minipage}[b]{\linewidth}\raggedright
可执行性
\end{minipage} \\
\midrule\noalign{}
\endhead
\bottomrule\noalign{}
\endlastfoot
1.19 & 小规模 LLM SFT 实验：选公开 instruction dataset (几千条)，LoRA
微调 0.5B-3B 模型 & P2 & 1.9 & 20h & {[}A{]}{[}LLM{]} & ⚡ \\
1.20 & 比较 SCX selection vs random/diversity/high-loss/uncertainty
(20\% data) & P2 & 1.19 & 10h & {[}A{]}{[}LLM{]} & ⏳ \\
1.21 & 写 LLM-SCX 的定位文档：SCX 不是通用 judge，而是''评价器编排框架''
& P2 & 1.19 & 4h & {[}LLM{]}{[}A{]} & ⚡ \\
\end{longtable}

\begin{center}\rule{0.5\linewidth}{0.5pt}\end{center}

\subsection{Phase 2: 中期（3-6
月）}\label{phase-2-ux4e2dux671f3-6-ux6708}

\subsubsection{论文深化}\label{ux8bbaux6587ux6df1ux5316}

\begin{longtable}[]{@{}
  >{\raggedright\arraybackslash}p{(\linewidth - 12\tabcolsep) * \real{0.0216}}
  >{\raggedright\arraybackslash}p{(\linewidth - 12\tabcolsep) * \real{0.6619}}
  >{\raggedright\arraybackslash}p{(\linewidth - 12\tabcolsep) * \real{0.0432}}
  >{\raggedright\arraybackslash}p{(\linewidth - 12\tabcolsep) * \real{0.0576}}
  >{\raggedright\arraybackslash}p{(\linewidth - 12\tabcolsep) * \real{0.0432}}
  >{\raggedright\arraybackslash}p{(\linewidth - 12\tabcolsep) * \real{0.0432}}
  >{\raggedright\arraybackslash}p{(\linewidth - 12\tabcolsep) * \real{0.1295}}@{}}
\toprule\noalign{}
\begin{minipage}[b]{\linewidth}\raggedright
\#
\end{minipage} & \begin{minipage}[b]{\linewidth}\raggedright
TODO
\end{minipage} & \begin{minipage}[b]{\linewidth}\raggedright
优先级
\end{minipage} & \begin{minipage}[b]{\linewidth}\raggedright
依赖
\end{minipage} & \begin{minipage}[b]{\linewidth}\raggedright
预估
\end{minipage} & \begin{minipage}[b]{\linewidth}\raggedright
维度
\end{minipage} & \begin{minipage}[b]{\linewidth}\raggedright
可执行性
\end{minipage} \\
\midrule\noalign{}
\endhead
\bottomrule\noalign{}
\endlastfoot
2.1 & SCX-Sim 大论文：多保真调度 + 科学工程仿真数据价值（DFT/MLIP +
FEM/PDE + OPC/CMP toy） & P2 & 1.5, 1.7 & 80h & {[}A{]}{[}C{]} & ⏳等
SCX-MLIP 投稿 \\
2.2 & 投稿目标：Nature Computational Science / Nature Communications /
Nature Machine Intelligence & P2 & 2.1 & 投稿日 & {[}A{]} & ⏳ \\
2.3 & FEM/PDE toy demo：悬臂梁 + 压电双悬臂梁 + Poisson/elasticity PDE &
P2 & 0.11 & 20h & {[}C{]}{[}A{]} & ⚡ \\
2.4 & SCX-Semi toy demo：简化光刻模型 + layout pattern state + 多 OPC
专家路由 & P2 & 0.11 & 25h & {[}C{]}{[}A{]} & ⚡ \\
2.5 & 将 SCX 的 identifiability limits 小节正式化（论文必要内容） & P2 &
1.7 & 6h & {[}A{]} & ⚡ \\
\end{longtable}

\subsubsection{商业拓展}\label{ux5546ux4e1aux62d3ux5c55}

\begin{longtable}[]{@{}
  >{\raggedright\arraybackslash}p{(\linewidth - 12\tabcolsep) * \real{0.0284}}
  >{\raggedright\arraybackslash}p{(\linewidth - 12\tabcolsep) * \real{0.5461}}
  >{\raggedright\arraybackslash}p{(\linewidth - 12\tabcolsep) * \real{0.0426}}
  >{\raggedright\arraybackslash}p{(\linewidth - 12\tabcolsep) * \real{0.1702}}
  >{\raggedright\arraybackslash}p{(\linewidth - 12\tabcolsep) * \real{0.1064}}
  >{\raggedright\arraybackslash}p{(\linewidth - 12\tabcolsep) * \real{0.0496}}
  >{\raggedright\arraybackslash}p{(\linewidth - 12\tabcolsep) * \real{0.0567}}@{}}
\toprule\noalign{}
\begin{minipage}[b]{\linewidth}\raggedright
\#
\end{minipage} & \begin{minipage}[b]{\linewidth}\raggedright
TODO
\end{minipage} & \begin{minipage}[b]{\linewidth}\raggedright
优先级
\end{minipage} & \begin{minipage}[b]{\linewidth}\raggedright
依赖
\end{minipage} & \begin{minipage}[b]{\linewidth}\raggedright
预估
\end{minipage} & \begin{minipage}[b]{\linewidth}\raggedright
维度
\end{minipage} & \begin{minipage}[b]{\linewidth}\raggedright
可执行性
\end{minipage} \\
\midrule\noalign{}
\endhead
\bottomrule\noalign{}
\endlastfoot
2.6 & 找 1-2 个 paid pilot 甲方（材料计算团队 / AI 数据公司 /
医学影像团队） & P2 & 1.15, 1.9 & 1-3 月 & {[}B{]} & ⏳ \\
2.7 & Pilot 交付物：数据价值审计报告（冗余分析 + 噪声风险 + 高价值状态 +
建议动作） & P2 & 2.6 & 每 pilot 2-4 周 & {[}B{]} & ⏳ \\
2.8 & 设计 SCX Consortium
产业联盟提案（会员费、参与规则、CLA、中立治理） & P3 & 2.6（需至少 2
家兴趣方） & 20h & {[}B{]}{[}IP{]} & ⏳ \\
2.9 & 建立 Contributor License Agreement (CLA) 模板 & P2 & 无 & 4h &
{[}IP{]} & ⚡ \\
2.10 & 设计三层商业产品线：SCX-Core(开源) / SCX-Health(开源) /
SCX-Pro(商业闭源) & P2 & 无 & 6h & {[}B{]}{[}C{]} & ⚡ \\
\end{longtable}

\subsubsection{SCX-Potential
Compiler（势函数编译器）}\label{scx-potential-compilerux52bfux51fdux6570ux7f16ux8bd1ux5668}

\begin{longtable}[]{@{}
  >{\raggedright\arraybackslash}p{(\linewidth - 12\tabcolsep) * \real{0.0339}}
  >{\raggedright\arraybackslash}p{(\linewidth - 12\tabcolsep) * \real{0.6864}}
  >{\raggedright\arraybackslash}p{(\linewidth - 12\tabcolsep) * \real{0.0508}}
  >{\raggedright\arraybackslash}p{(\linewidth - 12\tabcolsep) * \real{0.0763}}
  >{\raggedright\arraybackslash}p{(\linewidth - 12\tabcolsep) * \real{0.0339}}
  >{\raggedright\arraybackslash}p{(\linewidth - 12\tabcolsep) * \real{0.0508}}
  >{\raggedright\arraybackslash}p{(\linewidth - 12\tabcolsep) * \real{0.0678}}@{}}
\toprule\noalign{}
\begin{minipage}[b]{\linewidth}\raggedright
\#
\end{minipage} & \begin{minipage}[b]{\linewidth}\raggedright
TODO
\end{minipage} & \begin{minipage}[b]{\linewidth}\raggedright
优先级
\end{minipage} & \begin{minipage}[b]{\linewidth}\raggedright
依赖
\end{minipage} & \begin{minipage}[b]{\linewidth}\raggedright
预估
\end{minipage} & \begin{minipage}[b]{\linewidth}\raggedright
维度
\end{minipage} & \begin{minipage}[b]{\linewidth}\raggedright
可执行性
\end{minipage} \\
\midrule\noalign{}
\endhead
\bottomrule\noalign{}
\endlastfoot
2.11 & 构建 Expert Registry：登记 ACE/NEP/MACE/DeepMD 势函数的 metadata
& P2 & 0.13 & 10h & {[}C{]}{[}A{]} & ⚡ \\
2.12 & 实现 energy gauge alignment（species shift + coefficient gauge
fixing） & P2 & 0.13 & 15h & {[}C{]}{[}A{]} & ⚡ \\
2.13 & 构建 ACE-State Encoder / ACE-Critic（统一状态表示） & P2 & 0.13 &
10h & {[}C{]}{[}A{]} & ⚡ \\
2.14 & 构建 Expert Reliability Map（每个势函数在每个状态下的估计误差） &
P2 & 2.13 & 10h & {[}C{]}{[}A{]} & ⏳ \\
2.15 & 构建 Input Quality Judge（OOD 检测 + 专家分歧 + 误差风险） & P2 &
2.13 & 10h & {[}C{]}{[}A{]} & ⏳ \\
2.16 & 最小可行性验证：ACE + NEP + MACE + DeepMD 在 AlN/GaN 上做状态条件
ensemble 和蒸馏 & P2 & 2.11-2.15 & 30h & {[}C{]}{[}A{]} & ⏳ \\
2.17 & 验证：SCX teacher \textgreater{} naive ensemble \textgreater{}
best single expert & P2 & 2.16 & 10h & {[}A{]} & ⏳ \\
2.18 & 验证：SCX distillation student 保持 endpoint 不遗忘、OOD 风险更低
& P2 & 2.16 & 10h & {[}A{]} & ⏳ \\
\end{longtable}

\subsubsection{代码平台化}\label{ux4ee3ux7801ux5e73ux53f0ux5316}

\begin{longtable}[]{@{}
  >{\raggedright\arraybackslash}p{(\linewidth - 12\tabcolsep) * \real{0.0328}}
  >{\raggedright\arraybackslash}p{(\linewidth - 12\tabcolsep) * \real{0.7377}}
  >{\raggedright\arraybackslash}p{(\linewidth - 12\tabcolsep) * \real{0.0492}}
  >{\raggedright\arraybackslash}p{(\linewidth - 12\tabcolsep) * \real{0.0328}}
  >{\raggedright\arraybackslash}p{(\linewidth - 12\tabcolsep) * \real{0.0328}}
  >{\raggedright\arraybackslash}p{(\linewidth - 12\tabcolsep) * \real{0.0492}}
  >{\raggedright\arraybackslash}p{(\linewidth - 12\tabcolsep) * \real{0.0656}}@{}}
\toprule\noalign{}
\begin{minipage}[b]{\linewidth}\raggedright
\#
\end{minipage} & \begin{minipage}[b]{\linewidth}\raggedright
TODO
\end{minipage} & \begin{minipage}[b]{\linewidth}\raggedright
优先级
\end{minipage} & \begin{minipage}[b]{\linewidth}\raggedright
依赖
\end{minipage} & \begin{minipage}[b]{\linewidth}\raggedright
预估
\end{minipage} & \begin{minipage}[b]{\linewidth}\raggedright
维度
\end{minipage} & \begin{minipage}[b]{\linewidth}\raggedright
可执行性
\end{minipage} \\
\midrule\noalign{}
\endhead
\bottomrule\noalign{}
\endlastfoot
2.19 & 设计 SCX-Core
抽象接口：\texttt{SCXStateEncoder}（featurize/distance/cluster/state\_statistics）
& P2 & 0.11 & 8h & {[}C{]} & ⚡ \\
2.20 & 实现 SCX-MLIP Encoder（ACE descriptor） & P2 & 2.19 & 6h &
{[}C{]} & ⚡ \\
2.21 & 实现 SCX-Vision Encoder（DINO/CLIP/ViT embedding） & P2 & 2.19 &
6h & {[}C{]} & ⚡ \\
2.22 & 实现 SCX-FEM Encoder（几何/网格/响应 embedding） & P2 & 2.19 & 8h
& {[}C{]} & ⚡ \\
2.23 & 写数据价值策略模块：per-state budget allocation & P2 & 2.19 & 6h
& {[}C{]}{[}A{]} & ⚡ \\
2.24 & 写训练集构造策略模块：每个状态保留多少数据、加权/采样策略 & P2 &
2.19 & 6h & {[}C{]}{[}A{]} & ⚡ \\
2.25 & 写模型选择模块：根据数据状态推荐 NEP/MACE/ACE/DeepMD 等后端 & P2
& 2.19 & 6h & {[}C{]}{[}A{]} & ⚡ \\
\end{longtable}

\subsubsection{可穿戴健康数据}\label{ux53efux7a7fux6234ux5065ux5eb7ux6570ux636e}

\begin{longtable}[]{@{}
  >{\raggedright\arraybackslash}p{(\linewidth - 12\tabcolsep) * \real{0.0388}}
  >{\raggedright\arraybackslash}p{(\linewidth - 12\tabcolsep) * \real{0.6893}}
  >{\raggedright\arraybackslash}p{(\linewidth - 12\tabcolsep) * \real{0.0583}}
  >{\raggedright\arraybackslash}p{(\linewidth - 12\tabcolsep) * \real{0.0388}}
  >{\raggedright\arraybackslash}p{(\linewidth - 12\tabcolsep) * \real{0.0388}}
  >{\raggedright\arraybackslash}p{(\linewidth - 12\tabcolsep) * \real{0.0583}}
  >{\raggedright\arraybackslash}p{(\linewidth - 12\tabcolsep) * \real{0.0777}}@{}}
\toprule\noalign{}
\begin{minipage}[b]{\linewidth}\raggedright
\#
\end{minipage} & \begin{minipage}[b]{\linewidth}\raggedright
TODO
\end{minipage} & \begin{minipage}[b]{\linewidth}\raggedright
优先级
\end{minipage} & \begin{minipage}[b]{\linewidth}\raggedright
依赖
\end{minipage} & \begin{minipage}[b]{\linewidth}\raggedright
预估
\end{minipage} & \begin{minipage}[b]{\linewidth}\raggedright
维度
\end{minipage} & \begin{minipage}[b]{\linewidth}\raggedright
可执行性
\end{minipage} \\
\midrule\noalign{}
\endhead
\bottomrule\noalign{}
\endlastfoot
2.26 & 公开可穿戴数据 toy benchmark（心率/活动识别/睡眠状态/运动传感器）
& P2 & 1.9 & 15h & {[}C{]}{[}A{]} & ⚡ \\
2.27 & SCX-Health
增加可穿戴模块：连续数据冗余压缩、运动伪影噪声、长尾异常状态 & P2 & 2.26
& 15h & {[}C{]}{[}A{]} & ⏳ \\
\end{longtable}

\begin{center}\rule{0.5\linewidth}{0.5pt}\end{center}

\subsection{Phase 3: 长期（6-12
月）}\label{phase-3-ux957fux671f6-12-ux6708}

\subsubsection{论文发表（最高目标）}\label{ux8bbaux6587ux53d1ux8868ux6700ux9ad8ux76eeux6807}

\begin{longtable}[]{@{}
  >{\raggedright\arraybackslash}p{(\linewidth - 12\tabcolsep) * \real{0.0200}}
  >{\raggedright\arraybackslash}p{(\linewidth - 12\tabcolsep) * \real{0.7067}}
  >{\raggedright\arraybackslash}p{(\linewidth - 12\tabcolsep) * \real{0.0400}}
  >{\raggedright\arraybackslash}p{(\linewidth - 12\tabcolsep) * \real{0.1133}}
  >{\raggedright\arraybackslash}p{(\linewidth - 12\tabcolsep) * \real{0.0400}}
  >{\raggedright\arraybackslash}p{(\linewidth - 12\tabcolsep) * \real{0.0267}}
  >{\raggedright\arraybackslash}p{(\linewidth - 12\tabcolsep) * \real{0.0533}}@{}}
\toprule\noalign{}
\begin{minipage}[b]{\linewidth}\raggedright
\#
\end{minipage} & \begin{minipage}[b]{\linewidth}\raggedright
TODO
\end{minipage} & \begin{minipage}[b]{\linewidth}\raggedright
优先级
\end{minipage} & \begin{minipage}[b]{\linewidth}\raggedright
依赖
\end{minipage} & \begin{minipage}[b]{\linewidth}\raggedright
预估
\end{minipage} & \begin{minipage}[b]{\linewidth}\raggedright
维度
\end{minipage} & \begin{minipage}[b]{\linewidth}\raggedright
可执行性
\end{minipage} \\
\midrule\noalign{}
\endhead
\bottomrule\noalign{}
\endlastfoot
3.1 & SCX-Sim 论文正式投稿（Nature Computational Science / Nature
Communications / Nature Machine Intelligence） & P3 & 2.1-2.5 & 投稿日 &
{[}A{]} & ⏳ \\
3.2 & SCX-Health 论文投稿（npj Digital Medicine / Medical Image Analysis
/ Nature Communications 视数据强度） & P3 & 1.14 & 投稿日 & {[}A{]} &
⏳ \\
3.3 & SCX-Potential Compiler 论文（目标 npj Computational Materials /
Nature Communications） & P3 & 2.16-2.18 & 40h & {[}A{]} & ⏳ \\
3.4 & 如医学数据很强（多中心 + 医生标注），升级 SCX-Health 投稿 Nature
Biomedical Engineering & P3 & 3.2（需更硬数据） & 投稿日 & {[}A{]} &
⏳ \\
\end{longtable}

\subsubsection{商业落地}\label{ux5546ux4e1aux843dux5730}

\begin{longtable}[]{@{}
  >{\raggedright\arraybackslash}p{(\linewidth - 12\tabcolsep) * \real{0.0267}}
  >{\raggedright\arraybackslash}p{(\linewidth - 12\tabcolsep) * \real{0.6467}}
  >{\raggedright\arraybackslash}p{(\linewidth - 12\tabcolsep) * \real{0.0400}}
  >{\raggedright\arraybackslash}p{(\linewidth - 12\tabcolsep) * \real{0.1467}}
  >{\raggedright\arraybackslash}p{(\linewidth - 12\tabcolsep) * \real{0.0400}}
  >{\raggedright\arraybackslash}p{(\linewidth - 12\tabcolsep) * \real{0.0467}}
  >{\raggedright\arraybackslash}p{(\linewidth - 12\tabcolsep) * \real{0.0533}}@{}}
\toprule\noalign{}
\begin{minipage}[b]{\linewidth}\raggedright
\#
\end{minipage} & \begin{minipage}[b]{\linewidth}\raggedright
TODO
\end{minipage} & \begin{minipage}[b]{\linewidth}\raggedright
优先级
\end{minipage} & \begin{minipage}[b]{\linewidth}\raggedright
依赖
\end{minipage} & \begin{minipage}[b]{\linewidth}\raggedright
预估
\end{minipage} & \begin{minipage}[b]{\linewidth}\raggedright
维度
\end{minipage} & \begin{minipage}[b]{\linewidth}\raggedright
可执行性
\end{minipage} \\
\midrule\noalign{}
\endhead
\bottomrule\noalign{}
\endlastfoot
3.5 & 形成 SCX 数据价值审计标准化报告模板（8 个模块） & P3 & 2.6, 2.7 &
15h & {[}B{]} & ⏳ \\
3.6 & 建 SCX 认证体系（SCX-certified dataset / SCX-compressed dataset /
SCX expert reliability report） & P3 & 3.5 & 20h & {[}B{]}{[}IP{]} &
⏳ \\
3.7 & 找 CEO/COO 运营合伙人（你当 scientific founder） & P3 &
2.6（有成熟 pilot 后） & 1-3 月 & {[}B{]} & ⏳ \\
3.8 & 设 SCX Institute / SCX Lab（中立数据价值评估实验室） & P3 & 3.5,
3.7 & 数月 & {[}B{]}{[}IP{]} & ⏳ \\
3.9 & 找华为/小米健康团队做私有试点合作 & P3 & 2.26, 2.27 & 1-3 月 &
{[}B{]} & ⏳ \\
3.10 & 找 EDA/半导体公司（华大九天/概伦/广立微等）做 SCX-Semi 技术合作 &
P3 & 2.4 & 2-4 月 & {[}B{]} & ⏳ \\
\end{longtable}

\subsubsection{SCX 生态建设}\label{scx-ux751fux6001ux5efaux8bbe}

\begin{longtable}[]{@{}
  >{\raggedright\arraybackslash}p{(\linewidth - 12\tabcolsep) * \real{0.0348}}
  >{\raggedright\arraybackslash}p{(\linewidth - 12\tabcolsep) * \real{0.6348}}
  >{\raggedright\arraybackslash}p{(\linewidth - 12\tabcolsep) * \real{0.0522}}
  >{\raggedright\arraybackslash}p{(\linewidth - 12\tabcolsep) * \real{0.0870}}
  >{\raggedright\arraybackslash}p{(\linewidth - 12\tabcolsep) * \real{0.0522}}
  >{\raggedright\arraybackslash}p{(\linewidth - 12\tabcolsep) * \real{0.0696}}
  >{\raggedright\arraybackslash}p{(\linewidth - 12\tabcolsep) * \real{0.0696}}@{}}
\toprule\noalign{}
\begin{minipage}[b]{\linewidth}\raggedright
\#
\end{minipage} & \begin{minipage}[b]{\linewidth}\raggedright
TODO
\end{minipage} & \begin{minipage}[b]{\linewidth}\raggedright
优先级
\end{minipage} & \begin{minipage}[b]{\linewidth}\raggedright
依赖
\end{minipage} & \begin{minipage}[b]{\linewidth}\raggedright
预估
\end{minipage} & \begin{minipage}[b]{\linewidth}\raggedright
维度
\end{minipage} & \begin{minipage}[b]{\linewidth}\raggedright
可执行性
\end{minipage} \\
\midrule\noalign{}
\endhead
\bottomrule\noalign{}
\endlastfoot
3.11 & SCX-Health 开源项目建立社区（GitHub Stars, 文档, CI, 贡献指南） &
P3 & 2.26, 2.27 & 持续 & {[}C{]}{[}B{]} & ⏳ \\
3.12 & SCX-Semi 扩展：从 toy demo 升级到真实 PDK/layout
数据（需找到产业合作方） & P3 & 3.10 & 2-4 月 & {[}C{]}{[}A{]} & ⏳ \\
3.13 & SCX-Music 趣味 demo（作为展示 SCX 泛化能力的 showcase，非主线） &
P3 & 2.19 & 20h & {[}C{]}{[}B{]} & ⚡ \\
3.14 & LLM-SCX 方向：状态条件评价器编排框架（讨论/外延，不作为主战场） &
P3 & 1.19-1.21 & 30h & {[}LLM{]}{[}A{]} & ⏳ \\
\end{longtable}

\begin{center}\rule{0.5\linewidth}{0.5pt}\end{center}

\subsection{v4 通用模块化路线（2026-06-26
追加）}\label{v4-ux901aux7528ux6a21ux5757ux5316ux8defux7ebf2026-06-26-ux8ffdux52a0}

\begin{quote}
通用模块化架构是 SCX v4.0 的核心重构方向。详细规划见
\texttt{SCX\_规划\_v4\_通用模块化.md}，详细 TODO 见
\texttt{SCX\_TODO\_v4\_通用模块化.md}。
\end{quote}

\subsubsection{核心命题}\label{ux6838ux5fc3ux547dux9898}

\begin{quote}
``框架还是需要更通用一些，避免一直加模块迭代麻烦，应该模块化，不同的行业可以调用相同的代码。''
\end{quote}

\subsubsection{总体路线}\label{ux603bux4f53ux8defux7ebf}

\begin{verbatim}
v0.3.0 (当前)        v0.4.0 (Phase A)        v0.5.0 (Phase B)        v0.6.0 (Phase C/D)
──────────           ──────────────           ──────────────           ──────────────

紧耦合              抽象接口 + Encoder        YAML 驱动               插件系统
各领域独立pipeline  核心模块解耦              CLI 接口                 社区贡献
配置散落代码         3 领域示例               自动报告                  Benchmark
                    336 tests 通过           encoder 开发指南
\end{verbatim}

\subsubsection{Phase A: 核心重构 (2-4 周) 🔴
P0}\label{phase-a-ux6838ux5fc3ux91cdux6784-2-4-ux5468-p0}

\begin{longtable}[]{@{}
  >{\raggedright\arraybackslash}p{(\linewidth - 6\tabcolsep) * \real{0.1304}}
  >{\raggedright\arraybackslash}p{(\linewidth - 6\tabcolsep) * \real{0.2609}}
  >{\raggedright\arraybackslash}p{(\linewidth - 6\tabcolsep) * \real{0.3478}}
  >{\raggedright\arraybackslash}p{(\linewidth - 6\tabcolsep) * \real{0.2609}}@{}}
\toprule\noalign{}
\begin{minipage}[b]{\linewidth}\raggedright
\#
\end{minipage} & \begin{minipage}[b]{\linewidth}\raggedright
TODO
\end{minipage} & \begin{minipage}[b]{\linewidth}\raggedright
优先级
\end{minipage} & \begin{minipage}[b]{\linewidth}\raggedright
预估
\end{minipage} \\
\midrule\noalign{}
\endhead
\bottomrule\noalign{}
\endlastfoot
A.1 & 定义 \texttt{SCXStateEncoder} / \texttt{SCXExpert} /
\texttt{SCXDataPoint} 抽象基类 & P0 & 3d \\
A.2 & 提取 MLIP/Vision/CIFAR 三个 encoder 到 \texttt{encoders/} 目录 &
P0 & 3d \\
A.3 & 重构 \texttt{state/expert/valuation/action/online} 只依赖抽象接口
& P0 & 5d \\
A.4 & 创建 YAML 配置原型 + 3 领域示例配置 & P0 & 2d \\
A.5 & 验证：336 tests 全部通过，无回归 & P0 & 2d \\
\end{longtable}

\subsubsection{Phase B: 声明式配置 (2-3 周) 🟡
P1}\label{phase-b-ux58f0ux660eux5f0fux914dux7f6e-2-3-ux5468-p1}

\begin{longtable}[]{@{}llll@{}}
\toprule\noalign{}
\# & TODO & 优先级 & 预估 \\
\midrule\noalign{}
\endhead
\bottomrule\noalign{}
\endlastfoot
B.1 & YAML → SCXFramework 自动构建 & P1 & 6d \\
B.2 & \texttt{scx\ run\ -\/-config\ domains/mlip.yaml} CLI & P1 & 2d \\
B.3 & 审计报告自动生成 & P1 & 3d \\
B.4 & Encoder 开发指南 + YAML 配置参考 & P1 & 2d \\
\end{longtable}

\subsubsection{Phase C: 新领域扩展 (4-8 周) 🟡
P1}\label{phase-c-ux65b0ux9886ux57dfux6269ux5c55-4-8-ux5468-p1}

\begin{longtable}[]{@{}
  >{\raggedright\arraybackslash}p{(\linewidth - 6\tabcolsep) * \real{0.1304}}
  >{\raggedright\arraybackslash}p{(\linewidth - 6\tabcolsep) * \real{0.2609}}
  >{\raggedright\arraybackslash}p{(\linewidth - 6\tabcolsep) * \real{0.3478}}
  >{\raggedright\arraybackslash}p{(\linewidth - 6\tabcolsep) * \real{0.2609}}@{}}
\toprule\noalign{}
\begin{minipage}[b]{\linewidth}\raggedright
\#
\end{minipage} & \begin{minipage}[b]{\linewidth}\raggedright
TODO
\end{minipage} & \begin{minipage}[b]{\linewidth}\raggedright
优先级
\end{minipage} & \begin{minipage}[b]{\linewidth}\raggedright
预估
\end{minipage} \\
\midrule\noalign{}
\endhead
\bottomrule\noalign{}
\endlastfoot
C.1 & SCX-Robot: TrajectoryEncoder + robot.yaml + 失败/冗余实验 & P1 &
2-3w \\
C.2 & SCX-FEM: SimulationEncoder + fem.yaml + 多保真调度 & P1 & 2-3w \\
C.3 & SCX-LLM: TextEncoder + llm.yaml (minimal, 聚焦可验证任务) & P1 &
1-2w \\
\end{longtable}

\subsubsection{Phase D: 插件系统 (长期) 🟢
P2}\label{phase-d-ux63d2ux4ef6ux7cfbux7edf-ux957fux671f-p2}

\begin{longtable}[]{@{}llll@{}}
\toprule\noalign{}
\# & TODO & 优先级 & 预估 \\
\midrule\noalign{}
\endhead
\bottomrule\noalign{}
\endlastfoot
D.1 & 插件注册机制 + compress/online/influence 插件化 & P2 & 2w \\
D.2 & 社区贡献指南 + CLA & P2 & 1w \\
D.3 & SCX-Bench 标准化 API + 5-10 数据集 & P2 & 2w \\
\end{longtable}

\subsubsection{商业模式映射}\label{ux5546ux4e1aux6a21ux5f0fux6620ux5c04}

\begin{longtable}[]{@{}lll@{}}
\toprule\noalign{}
行业 & 产品 & 定价锚点 \\
\midrule\noalign{}
\endhead
\bottomrule\noalign{}
\endlastfoot
MLIP/DFT & SCX-Potential Compiler & 节省核时 \\
医学图像 & SCX-Health Audit & 节省标注成本 \\
机器人 & SCX-Robot Data Audit & 节省遥操作/真机时间 \\
工程仿真 & SCX-FEM Scheduler & 节省仿真时间 \\
工业质检 & SCX-Vision Audit & 减少人工复检量 \\
3D 打印 & SCX-AM Audit & 减少打样次数 \\
半导体 & SCX-Semi OPC Router & 节省 EDA license 费用 \\
\end{longtable}

\subsubsection{与旧路线的衔接}\label{ux4e0eux65e7ux8defux7ebfux7684ux8854ux63a5}

\begin{itemize}
\tightlist
\item
  \textbf{GPU 等待项}（CIFAR/MedMNIST/HAM10000 GPU 重跑）---
  保留，8月中旬
\item
  \textbf{DFT 等待项}（GaN/AlN VASP 结果）--- 保留，2-4周
\item
  \textbf{论文写作}（EGP Paper 1, SCX-Theory arXiv）--- 保留，并行推进
\item
  通用架构重构\textbf{不影响}论文进度，反而让实验更易复现
\end{itemize}

\begin{center}\rule{0.5\linewidth}{0.5pt}\end{center}

\subsection{依赖关系总图}\label{ux4f9dux8d56ux5173ux7cfbux603bux56fe}

\begin{verbatim}
Phase 0                              Phase 1                          Phase 2                              Phase 3
───────                              ───────                          ───────                              ───────

0.1 私人仓库 ─────────────────────────→ 1.4 时间戳存证
0.2 DEVELOPMENT_LOG.md                1.1 校外 IP 律师 ──────────────→ 1.2 专利价值判断 ─────────→ 1.3 专利申请
0.3 代码来源表
0.4 协议检查 ─────────────────────────→ 1.1
0.5 Clean-room 重写 ─────────────────→ 0.11 整理 v0.2.0 ────────────→ 2.19 SCX-Core 接口 ──────→ 2.20-2.25 各 Encoder
                                      │                                    │
0.9 EGP Paper 1 (ACE merging) ───────→ 1.5 SCX-MLIP 论文 ────────────→ 2.1 SCX-Sim 大论文 ─────→ 3.1 投稿 Nature 系列
0.10 SCX-Theory 草稿 ───────────────→ 1.7 arXiv 占坑 ───────────────→ 2.5 identifiability 小节
                                      │
0.13 ACE pipeline 整理 ──────────────→ 2.11-2.18 Potential Compiler ──→ 3.3 势函数编译器论文
                                      │
0.14 VASP 任务检查 ─────────────────→ 0.15 提交 ────────────────────→ 0.16 监控脚本
                                      │
                                      │
                                      1.9 SCX-Health 仓库 ──────────→ 1.10-1.13 医学实验 ────→ 1.14 SCX-Health 论文 ──→ 3.2 投稿
                                      │                                │
                                      1.15 One-pager ─────────────────→ 2.6 找 paid pilot ─────→ 2.7 交付审计报告 ───→ 3.5 标准化报告
                                                                        │                        │
                                      1.18 国产 EDA/TCAD 调研 ────────→ 2.4 SCX-Semi toy ───────→ 3.10 产业合作
                                                                                                3.7 找 CEO/COO ────→ 3.8 SCX Institute
\end{verbatim}

\begin{center}\rule{0.5\linewidth}{0.5pt}\end{center}

\subsection{核心风险提示}\label{ux6838ux5fc3ux98ceux9669ux63d0ux793a}

\begin{longtable}[]{@{}
  >{\raggedright\arraybackslash}p{(\linewidth - 4\tabcolsep) * \real{0.2800}}
  >{\raggedright\arraybackslash}p{(\linewidth - 4\tabcolsep) * \real{0.0800}}
  >{\raggedright\arraybackslash}p{(\linewidth - 4\tabcolsep) * \real{0.6400}}@{}}
\toprule\noalign{}
\begin{minipage}[b]{\linewidth}\raggedright
风险
\end{minipage} & \begin{minipage}[b]{\linewidth}\raggedright
严重程度
\end{minipage} & \begin{minipage}[b]{\linewidth}\raggedright
缓解措施
\end{minipage} \\
\midrule\noalign{}
\endhead
\bottomrule\noalign{}
\endlastfoot
学校主张 SCX IP 权属 & 🔴 高 & Phase 0 证据隔离 + 1.1 校外律师评估 +
切割 DFT/AlGaN 和 SCX 数学 \\
导师/同事抢先发表类似概念 & 🔴 高 & 1.7 arXiv 尽快占坑 + 0.2
开发日志存证 \\
大厂平替通用 LLM 数据评价 & 🟡 中 & 主战场不放在 LLM，放
DFT/MLIP/FEM/医学 \\
审稿人质疑 salami slicing & 🟡 中 &
每篇论文有独立贡献（见论文谱系设计） \\
论文实验不充分被拒 & 🟡 中 & 先出 MLIP 根据地论文，再扩展 Sim/Health \\
商业化过早/被甲方绑定 & 🟡 中 & 坚持非独占授权、CLA、不卖断 \\
VASP 任务占满 240 核影响他人 & 🟢 低 & 分批提交、监控负载 \\
医学数据伦理/隐私违规 & 🔴 高 &
先用公开数据，不碰真实患者数据；合规前置 \\
\end{longtable}

\begin{center}\rule{0.5\linewidth}{0.5pt}\end{center}

\subsection{五个维度的核心策略}\label{ux4e94ux4e2aux7ef4ux5ea6ux7684ux6838ux5fc3ux7b56ux7565}

\subsubsection{学术 {[}A{]}}\label{ux5b66ux672f-a}

\begin{verbatim}
EGP Paper 1 — ACE gauge merging (根据地, npj Computational Materials / PRM / JCTC)
  → SCX-Theory — 数学地基 (命名权, arXiv / TMLR / SIMODS / JMLR)
    → SCX-MLIP — 理论→应用闭环 (Nature Communications / npj Computational Materials)
      → SCX-Sim — 跨领域仿真 (Nature Computational Science)
        → SCX-Health — 医学数据估值 (npj Digital Medicine)
\end{verbatim}

总纲：5 篇论文谱系，EGP Paper 1 优先（有完整 DFT 数据），SCX
理论→应用→扩展。

\subsubsection{代码 {[}C{]}}\label{ux4ee3ux7801-c}

\begin{verbatim}
SCX-Core (抽象接口)
  ├── SCX-MLIP Encoder (ACE descriptor)
  ├── SCX-Vision Encoder (DINO/CLIP)
  ├── SCX-FEM Encoder (几何/网格)
  ├── SCX-Semi Encoder (layout/process)
  └── SCX-LLM Encoder (task embedding)
\end{verbatim}

开源策略：Core 开源，Encoders 按领域开源，Pro 版闭源。

\subsubsection{商业 {[}B{]}}\label{ux5546ux4e1a-b}

\begin{verbatim}
Phase 1: One-pager + 合作调研（零成本）
Phase 2: Paid pilots + 数据价值审计报告
Phase 3: 标准化报告 + 认证体系 + 产业联盟
Phase 4: SCX Institute (中立标准方) + CEO/COO 运营
\end{verbatim}

不卖断、不独占、保持中立。

\subsubsection{IP {[}IP{]}}\label{ip-ip}

\begin{verbatim}
证据隔离 (Phase 0)
  → 校外律师评估 (Phase 1)
    → 专利/软著/存证 (Phase 1-2)
      → CLA + 授权协议 (Phase 2)
        → 认证体系商标 (Phase 3)
\end{verbatim}

核心原则：先保护再公开，学校资产和个人资产切割。

\subsubsection{LLM {[}LLM{]}}\label{llm-llm}

\begin{verbatim}
最小验证 (Phase 1-2)
  → 方向定位：评价器编排框架，不是万能 judge
    → 作为 SCX 远期外延，不作为主战场
\end{verbatim}

明确：大厂内部化严重，SCX 不适合正面切入。

\begin{center}\rule{0.5\linewidth}{0.5pt}\end{center}

\subsection{总结：现在立刻做的 5
件事}\label{ux603bux7ed3ux73b0ux5728ux7acbux523bux505aux7684-5-ux4ef6ux4e8b}

\begin{enumerate}
\def\labelenumi{\arabic{enumi}.}
\tightlist
\item
  \textbf{私人仓库 + 证据链} (0.1, 0.2, 0.3) ------ IP 安全的生命线
\item
  \textbf{查协议} (0.4) ------ 了解权属风险范围
\item
  \textbf{Clean-room 重写 SCX 代码} (0.5, 0.11) ------ 从源头切割
\item
  \textbf{写 EGP Paper 1 草稿 + SCX-Theory arXiv 草稿} (0.9, 0.10)
  ------ 抢学术优先权
\item
  \textbf{建立 SCX-Health 开源仓库 + 跑 MedMNIST 实验} (1.9, 1.10)
  ------ 开源影响力入口
\end{enumerate}

\begin{quote}
\textbf{一句话战略：论文抢概念 + 轻量开源建生态 + 闭源工程炼商业 + IP
隔离保安全。大模型是远期外延，不是你的主战场。}
\end{quote}

\end{document}
